\documentclass[11pt,a4paper]{article}
\usepackage[margin=23mm]{geometry}
\usepackage{amsmath,amssymb,amsthm,booktabs,array,longtable}
\input{glyphtounicode}
\pdfgentounicode=1
\usepackage[hidelinks]{hyperref}
\hypersetup{hypertexnames=false}
\setlength{\parindent}{0pt}
\setlength{\parskip}{5pt}
\newcommand{\pos}[1]{[#1]_+}
\newcommand{\dd}{\,\mathrm d}
\newcommand{\Bud}{\mathcal B}
\newcommand{\He}{\mathcal H}
\newtheorem{proposition}{Proposition}
\newtheorem{theorem}{Theorem}
\begin{document}
\begin{center}
{\Large Standalone Step 3: information, actions and complete gradients}\\[4pt]
19 October 2026 --- selected H1/T1, S-E-P, R-E, O-R and W-F law
\end{center}

\textbf{Scope.} The user makes SUMO optional. The defined conditional model
is the present verification target. This note changes no equation, action,
observation, queue guard, renewal mass or completion boundary. It derives
new analytical comparison and sensitivity references; it does not run a
traffic trajectory, AD, numerical differences or an optimizer. Physical
admission and external usefulness are unestablished. Full equations are in
the \href{vsl_feedback8_full_mathematics.pdf}{18-page consolidated model}.

\section{What information is missing, and what closes prediction?}
Let $v\in[v_-,v_+]$, $L=L_p+b$, and let the initial homogeneous footprint
be $[0,\ell]$. Densities are $e_0,a_0=b_0$, total density $\rho$, mass
$M=\ell\rho$, and $r=e_0+a_0$. Physical parameters are common and known;
the initial terminal and recovery histories are empty. With $K=wk_j$,
\begin{equation}
 k_c(v)=\frac{K}{w+v},\quad C(v)=\frac{Kv}{w+v},\quad
 U(g)=\frac{wg}{K-g},\quad
 e_z=-h^{-1}\min\{e,(k_c(v)-r+e)/\gamma\}.
 \label{eq:h1}
\end{equation}
The admitted domain includes $v_+r<R\leq C(v_-)$, finite storage
$M<k_jb$, the original lane-density/far-lane bounds, and
$c_0=1+\delta\geq1$, $\beta\geq0$,
$c_0-\beta(1+v_+/w)>0$. The fixed coefficient is
$\chi=\beta/[C(v_+)v_+]$; it is not differentiated with advice.

The earlier H1 loading construction remains valid because all predecision
traffic stays upstream of the cutoff; the amended terminal is empty.
For common past speed $u$, $T=\ell/u$, prescribe on $[-T,0]$
\begin{equation}
 e_i^{\rm in}(t)=e_{0,i}e^{-ut/h},\quad
 a_i^{\rm in}(t)=r_i-e_i^{\rm in}(t),\quad b_i^{\rm in}=b_{0,i}.
 \label{eq:loading}
\end{equation}
Under the previously checked donor/loading bounds, both preparations have
the same occupied front $u(t+T)$ and identical lane-blind fields
\begin{equation}
 k_\Sigma=\rho\mathbf1_{0<x<u(t+T)},\quad
 q_\Sigma=uk_\Sigma,\quad
 N_\Sigma=\rho[u(t+T)-x]_+.
 \label{eq:collision}
\end{equation}
Thus even their full aggregate density, speed, longitudinal-flow and
cumulative-count histories cannot identify the actual initial ending mass.
Lane or lateral-event observations are extra channels, not part of this
collision. A lawful aggregate inference gets every admitted channel.

Given $E_0$, however,
\begin{equation}
 e_0=E_0/\ell,\qquad a_0=b_0=(M-E_0)/(2\ell).
 \label{eq:reconstruct}
\end{equation}
The conditional unique forward law then predicts all preparation, arrivals,
queues, speed marks, recovery history and departures. Aggregate-plus-$E_0$
(A+E) and explicit lane state (G) therefore give identical $J(v)$ for every
admitted advice, and identical derivatives wherever defined. Recovery
history is predicted state, not an extra initial observation. This does
not prove $E_0$ is the only possible sufficient encoding.

\section{A retained reference family and an exact negative interval}
Reuse the earlier analytical H1/T1 pair, rather than S13's replaced clock:
\begin{equation}
\begin{gathered}
v_-=u,\quad v_+=2u,\quad w=2u,\quad k_j=3\kappa,\quad\gamma=2,\\
R=39\kappa u/20,\quad R_B=\kappa u,\quad b=\ell,\quad
L_p=\ell+h\log2,\quad 0<\ell/h<\log(4/3),\\
\mathcal L:(e_0,a_0,b_0)=(3\kappa/5,3\kappa/10,3\kappa/10),\\
\mathcal H:(e_0,a_0,b_0)=(7\kappa/10,\kappa/4,\kappa/4).
\end{gathered}
\label{eq:family}
\end{equation}
Both have $\rho=6\kappa/5$ and the identical history
\eqref{eq:collision}. All original domain bounds hold; in particular
$C(u)=2\kappa u$, $M/(k_jb)=2/5$ and
$R-v_+\max_i r_i=\kappa u/20>0$.

Write $E_f(v)=e(L_p-\ell;v)$. The earlier explicit preparation calculation
gives
\begin{equation}
 E_{f,\mathcal L}=3\kappa/10,\quad
 E_{f,\mathcal H}(v_+)=125\kappa/352,\quad
 E'_{f,\mathcal H}(v_+)=225\kappa/(7744u)>0.
 \label{eq:front}
\end{equation}
H1 comparison makes front work
$W_f(v)=v[r+\delta E_f(v)]$ increasing with advice. At $v_+$,
\begin{equation}
 \frac{W_{f,\mathcal L}}{\kappa u}=\frac95+\frac35\delta,
 \qquad
 \frac{W_{f,\mathcal H}}{\kappa u}=\frac{19}{10}+\frac{125}{176}\delta.
 \label{eq:frontwork}
\end{equation}
Consequently, for the nonzero interval
\begin{equation}
 0\leq\delta\leq44/625,
 \qquad J_{\mathcal L}(v)=J_{\mathcal H}(v)
 =\frac{M(L-\ell/2)}v.
 \label{eq:negative}
\end{equation}
Both streams stay free for all advice. Equal free marks make mismatch
zero and free E creates no recovery; equality at the threshold remains
free under the selected guard. This result permits positive $\beta$ and
finite positive recovery acceleration. Hidden E mass exists but is
unnecessary for the best advice in this interval. The larger predecessor
small-work classification is not asserted for the coupled law.

\section{Decision value in the selected law: a bounded comparison proof}
\subsection{The auxiliary fixed-work reference}
Delete mismatch and retained recovery only for a reference, denoted $J_i^0$.
H1, arrivals, storage, initial information, completion and $c_0$ are unchanged.
A is free because its offer is below $R$; E is a reflected fluid queue with
service $(R-\lambda_A)/c_0$. This is exactly the earlier H1/T1 server, not
a replacement proposed for the selected law.

For sufficiently large fixed finite $c_0$, the earlier complete-time
calculation under \eqref{eq:family} proves:
\begin{enumerate}
\item $J^0_{\mathcal L}$ has the unique optimum $v_+$;
\item a feasible small decrease from $v_+$ strictly improves
$J^0_{\mathcal H}$, whereas it strictly worsens $J^0_{\mathcal L}$;
\item the two continuous reference objectives have disjoint minimizing sets.
\end{enumerate}
The branch is ensured by $c_0>17/3$. More precisely, writing the earlier
exact expression as
\begin{equation}
 J_i^0=\frac{c_0\overline E_i(v)^2}{2R}
       +I_{0,i}(v)+I_{1,i}(v)/c_0,
 \label{eq:oldexpand}
\end{equation}
with $\overline E_i=hE_{f,i}(1-e^{-\ell/h})$,
$A_+=\overline E_{\mathcal H}(v_+)
\overline E'_{\mathcal H}(v_+)/R>0$, it is sufficient to take
\begin{equation}
 c_0>\max\{17/3,(|I'_{0,\mathcal H}(v_+)|+
                     |I'_{1,\mathcal H}(v_+)|)/A_+\}.
 \label{eq:c0range}
\end{equation}
All geometry is fixed first. This recalls the archived exact proof; it
does not select an experimental work factor or claim physical support.

\subsection{Uniformly bounding the actual mismatch and recovery effects}
Now restore the selected W-F/R-E/O-R law with $\beta>0$ and finite $a_r>0$.
Its isolated loss area at zero mark and a uniform account bound are
\begin{equation}
 u_R=U(R),\quad W_0=\frac{K}{w+u_R}\frac{u_R^2}{2a_r},\quad
 0\leq p(t)\leq\varepsilon_p:=\frac{R W_0}{mc_0}.
 \label{eq:smallloss}
\end{equation}
Indeed $g_E\leq R/c_0$ and the retained-window integral is a subset of
the nonnegative time integral, bounded by $(R/c_0)W_0/m$. This argument
retains the actual renewal rule and any moving cutoff.

Let $a_* =v_+\max_i r_i<R$. Take
\begin{equation}
 0<\beta\leq1,\quad c_0-\beta(1+v_+/w)>0,\qquad
 \varepsilon_p<\min\{(R-a_*)/2,R/2\}.
 \label{eq:smallconditions}
\end{equation}
Then $\Bud\geq R-\varepsilon_p>a_*\geq\lambda_A$: A remains free,
not by suppressing its queue, but because its actual offer always fits.
The selected E flow while queued lies between
\begin{equation}
 \mu_-(t)=\frac{R-\varepsilon_p-\lambda_A(t)}{c_0+\beta}
 \ \leq g_E(t)\leq
 \mu_0(t)=\frac{R-\lambda_A(t)}{c_0}.
 \label{eq:servicebracket}
\end{equation}
Here $\lambda_A=0$ outside its pulse. Both comparison services are positive.
If E is empty, its free continuation serves its offer or activates only
on a compatible strict shortfall. These guards preserve queue comparison.

For the same E arrivals, denote the reflected queues driven by
$\mu_0,\mu_-$ by $q_0,q_-$. Comparing derivatives at a first crossing of
positive queues, and using nonnegative reflection at zero, gives
\begin{equation}
 q_0(t)\leq q_E(t)\leq q_-(t),\qquad q_A(t)=0.
 \label{eq:queuebracket}
\end{equation}
This comparison does not assert that the actual state-dependent queue
itself is a fixed-capacity reflection. It accommodates the free/queued
speed-mark change at birth.

Moreover,
\begin{equation}
 0\leq\mu_0-\mu_-
 =\frac{\beta(R-\lambda_A)+c_0\varepsilon_p}{c_0(c_0+\beta)}
 \leq\varepsilon_\mu:=\frac{\varepsilon_p}{c_0}
                           +\frac{R\beta}{c_0^2}.
 \label{eq:servicedifference}
\end{equation}
The reflection representation as the supremum over net arrivals on
$[s,t]$ implies $0\leq q_--q_0\leq t\varepsilon_\mu$.
Choose a fixed $H>L/v_-+2(c_0+1)M/R$, independent of the small
parameters. All compared cohorts clear by $H$. Integrating the queue
bracket proves the new uniform objective estimate
\begin{equation}
 0\leq J_i(v)-J_i^0(v)\leq
 \varepsilon_J:=\frac{H^2}{2}\varepsilon_\mu
 \quad\hbox{for both preparations and every advice.}
 \label{eq:uniformvalue}
\end{equation}
Since $\varepsilon_p=O(a_r^{-1})$, the bound can be made arbitrarily
small with positive $\beta$ and finite positive $a_r$. These are
mathematical existence choices, not recommended traffic settings.

\subsection{A common aggregate action has unavoidable regret}
Define the positive reference separation
\begin{equation}
 \Gamma_0=\min_{v\in[v_-,v_+]}\max_{i\in\{\mathcal L,\mathcal H\}}
 \{J_i^0(v)-\min_zJ_i^0(z)\}>0.
 \label{eq:refregret}
\end{equation}
It is positive by continuity, compactness and disjoint minimizing sets.
For the actual objectives let $V_i=\inf_vJ_i(v)$, making no additional
global smoothness or attainment assumption. From \eqref{eq:uniformvalue},
\begin{equation}
 \inf_v\max_i\{J_i(v)-V_i\}\geq\Gamma_0-2\varepsilon_J>0
 \quad\hbox{if }\varepsilon_J<\Gamma_0/4.
 \label{eq:actualregret}
\end{equation}
Therefore every deterministic aggregate-only action has nonzero regret
for at least one preparation, even allowing optimal lawful inference from
the identical full history. A+E and G can distinguish the preparations.
The same conclusion concerns arbitrarily near-optimal actions if an
infimum is not attained; no general minimum-existence theorem is supplied.

Choose a fixed feasible $\Delta v>0$ from the strict reference slope
contrast. The two reference cost changes between $v_+$ and $v_+-\Delta v$
have opposite signs. Taking $2\varepsilon_J$ below both strict gaps
preserves those signs for the actual selected law. This establishes
opposite \emph{finite-perturbation} usefulness, not automatically opposite
infinitesimal gradients at $v_+$.

\textbf{Limits of this result.} The constructed range may need large work,
weak mismatch and fast recovery. A stays free there. The earlier complete
A-delay/active-renewal witness is a different range; no joint strong-delay
and action-split witness is claimed. Neither range is calibrated or a
default evaluation setting. The proof constructs existence transparently;
it does not license selecting only favorable evaluation cases.

\section{The full sensitivity path on regular event regimes}
Use $e_v(z;v)$ for the derivative at fixed physical progress. On each
smooth H1 branch its variational equation follows directly from
\eqref{eq:h1}; $e_v(0)=0$ and
\begin{equation}
 (e_v)_z=\begin{cases}
 -e_v/h,&\text{donor},\\
 -(e_v+k_c'(v))/(\gamma h),&\text{receiving},
 \end{cases}\qquad k_c'(v)=-K/(w+v)^2.
 \label{eq:prepderivative}
\end{equation}
The H1 vector fields match at their switch. Its moving location matters,
but there is no state-sensitivity jump from unequal vector fields there.

Each initial position $x$ becomes eligible at $(L-x)/v$. Exactly,
\begin{equation}
 P_E(t;v)=\int_0^\ell e(L_p-x;v)
             \mathbf1_{(L-x)/v\leq t}\dd x.
 \label{eq:arrival}
\end{equation}
For an interior pulse time, put $x_*=L-vt$, $z_*=vt-b$; then
\begin{align}
 P_{E,v}&=\int_{x_*}^\ell e_v(L_p-x;v)\dd x+t e(z_*;v),\nonumber\\
 \lambda_E&=v e(z_*;v),\qquad
 \lambda_{E,v}=e+v e_v+vt e_z,\nonumber\\
 \lambda_{A,v}&=r-\lambda_{E,v},\quad
 t_f'=-(L-\ell)/v^2,\quad t_r'=-L/v^2.
 \label{eq:arrivalderivative}
\end{align}
Before the front $P_{E,v}=0$; after the rear it is the full integral of
$e_v$. Front/rear changes are treated with their one-sided values. Holding
travel delays or arrival times fixed removes a real causal channel.

For queued E, $\He(g;a,u)=g[c_0+\chi a(u-U(g))_+]=x$. On positive
mismatch, its regular implicit derivative is
\begin{align}
 \He_g&=c_0+\chi a[u-U(g)-gU'(g)]\geq\eta>0,\nonumber\\
 \He_a&=\chi g[u-U(g)],\quad \He_u=\chi ag,\qquad
 g_v=\frac{x_v-\He_a a_v-\He_u u_v}{\He_g}.
 \label{eq:rootderivative}
\end{align}
On the inactive side, $\He_g=c_0$, $\He_a=\He_u=0$.
Free A uses $a=\lambda_A$, $u=v$, $x=\Bud-\lambda_A$.
Both queued use $a=x=\Bud/2$, $u=U(\Bud/2)$.
Thus there is no direct advice term in fully queued local work at fixed
budget; advice still changes the budget and arrivals. Differentiating an
advice-dependent normalizer would reintroduce the rejected W-A channel.

\section{Recovery differentiation must include events and retirement}
For fixed physical parameters, write
\begin{equation}
 p(t)=\frac1m\int_{s_m(t)}^t
 D(t-s;u_E(s))g_E(s)\dd s,
 \label{eq:recoverytime}
\end{equation}
where $s_m=0$ before fill and, after fill,
$Y_E(s_m)=Y_E(t)-m$. Zero-rate intervals carry no admitted mass. On a
regular filled-window chart the old cutoff lies at positive E rate and
away from a mark jump. Put $Z=\partial_vY_E$ and
$D_o=D(t-s_m;u_E(s_m))$. Differentiating the rank equation gives
\begin{equation}
 g_E(s_m)\partial_vs_m=Z(t)-Z(s_m).
 \label{eq:cutoffderivative}
\end{equation}
Let $\tau_k(v)$ be the finite source/mark switching times strictly inside
the retained time window, and set
$\Phi_k^\pm(t)=D(t-\tau_k;u_E(\tau_k^\pm))g_E(\tau_k^\pm)$.
Piecewise Leibniz differentiation yields
\begin{align}
 p_v(t)=\frac1m\bigg[&\int_{s_m}^t
       \{D_u u_{E,v}g_E+Dg_{E,v}\}\dd s
       +\sum_k(\Phi_k^- -\Phi_k^+)\tau_k'\nonumber\\
       &-\mathbf1_{Y_E(t)>m}D_o\{Z(t)-Z(s_m)\}\bigg],
 \qquad \Bud_v=-p_v.
 \label{eq:recoveryderivative}
\end{align}
Integral derivatives are within the smooth pieces; the sum accounts for
their moving endpoints. The final term is retirement of old rank weight.
It has no inverse old-rate factor after cancellation. Omitting it
differentiates a different, all-history source. Profile support cutoffs
add no jump term because $D$ is continuous and zero there. At a fill,
cutoff/mark encounter or rate plateau, use the appropriate event/one-sided
analysis; the displayed regular formula is not a universal differentiability
theorem at such boundaries.

The active isolated profile uses
$S=\sqrt{(w+u)^2+2wa_r\theta}$ and has
\begin{equation}
 D_u=-Kw(w+u)/S^3,\qquad D_\theta=-Kw^2a_r/S^3.
 \label{eq:kernelderivative}
\end{equation}
These are zero on its inactive side. The future update must retain the
speed/age history, rather than regard the current meter reading as all
dynamic state.

At a transverse continuous-state event with guard
$\psi(\tau,X(\tau^-;v),v)=0$, differentiate the guard and match the same
physical state at the shifted event time:
\begin{equation}
 \tau'=-\frac{\psi_v+\psi_X X_v^-}{\psi_t+\psi_X F^-},\qquad
 X_v^+=X_v^-+(F^- -F^+)\tau'.
 \label{eq:eventderivative}
\end{equation}
For an A birth the guard is its free-feasibility margin
$\Bud-\lambda_A$, not the identically zero free A queue. For a transverse
clearance it is the positive queue reaching zero. This is the
identity-reset event-timing correction; see Kong et al., Section III-A/B.
The history-dependent account still requires our separate
\eqref{eq:recoveryderivative}; a finite-dimensional saltation citation
does not prove that history reduction away.

\section{Complete TST, regular derivatives and nonregular cases}
For each fixed parameter case choose one action-independent horizon
\begin{equation}
 H_*>L/v_-+c_{\max}M/B_{\min},\qquad B_{\min}>0
 \label{eq:horizon}
\end{equation}
from the already proved uniform budget bound. All mass has crossed by it;
remaining recovery of already completed traffic need not be integrated.
The two equivalent objective and sensitivity expressions are
\begin{align}
 J&=\int_0^{H_*}(M-Y_A-Y_E-Y_B)\dd t
   =\frac{M(L-\ell/2)}v+\int_0^{H_*}(q_A+q_E)\dd t,\nonumber\\
 J'&=-\int_0^{H_*}(Z_A+Z_E+Z_B)\dd t
   =-\frac{M(L-\ell/2)}{v^2}
        +\int_0^{H_*}(q_{A,v}+q_{E,v})\dd t.
 \label{eq:tstderivative}
\end{align}
Queues are continuous and zero at their episode endpoints, so queue-area
endpoint values contribute zero. Their event-time effects on the subsequent
state, flows and history do not vanish. In particular, an A birth can
change service to E while E already has a positive queue.

\begin{proposition}[Regular-chart reference]
On a neighborhood with a fixed finite event sequence, transverse guards,
smooth branches and regular renewal cutoffs, the selected local causal
construction admits the sensitivities above. They determine the
complete-objective derivative, including the terminal clearance tail.
\end{proposition}
\begin{proof}
Within each history chart, the strictly regular root and piecewise smooth
Leibniz map have parameter derivatives. Differentiate the locally uniform
Volterra contraction used in the forward construction. This gives the
linear causal sensitivity equation. At each of the finitely many transverse
events, the implicit-function guard derivative and state matching give
\eqref{eq:eventderivative}; split history integrals at those moving times.
Induct through the fixed event list. Finite bounds on the chart justify
differentiating the fixed-horizon objective. This proves a local statement,
not regularity of every advice or event sequence.
\end{proof}

Forward-admitted grazing, coincident events, queue-activation thresholds,
renewal fill/cutoff encounters and root/profile kinks remain in reporting.
Check continuity and feasible one-sided difference behavior separately.
If no finite ordinary or one-sided derivative exists, label it so and use
actual finite-perturbation costs for that comparison. Do not smooth,
freeze events, delete boundary cases or reinterpret a jump quotient as a
correct gradient. No general continuity, global minimum attainment or
opposite infinitesimal slopes for the selected coupled objectives is
proved by the uniform value result in Section 3.

The literature basis is deliberately bounded. Kong et al. explains the
event-timing correction and its transversality limits; it is not a traffic
or infinite-history theorem. Delle Monache et al. establishes that VSL
cost variations and gradient control already have a literature precedent,
under a different LWR/outflow-tracking task; its equations are not our TST
gradient. The model-specific arrival, root, renewal and comparison
derivations here are our deductions from the accepted law.

\section{Verification contract and present evidence limit}
Later execution must distinguish the continuous analytical derivative,
the implemented forward model's directional differences and its AD
derivative. Agreement of the last two alone does not prove convergence to
the defined model. Use double precision, feasible step/refinement ladders,
event signatures, complete clearance and both forms of
\eqref{eq:tstderivative}. Validate the value/change of actual complete TST
against the predicted direction and a strong scalar-search reference;
information advantage does not imply AD or lane-state superiority.

This review checks symbolic identities and exact static inequalities only.
No numerical traffic/ODE trajectory, finite-difference gradient, AD,
calibration or policy optimization is performed. Numerical tolerances and
resource/solver settings belong to a separately approved execution plan.
The positive existence and negative interval are both retained; neither
chooses the main evaluation parameters or establishes physical prevalence.

\begin{thebibliography}{9}
\bibitem{h1ref} Archived own derivation.
\href{vsl_step3_information_and_sensitivity.pdf}{Step 3: information and full-time sensitivity},
14 October 2026, Sections 1--7. Used only as the explicitly deleted-mechanism
reference and predecision loading construction.
\bibitem{current} Archived own derivation.
\href{vsl_feedback8_full_mathematics.pdf}{Full current mathematics},
16 October 2026. Selected equations and conditional forward bounds.
\bibitem{kong} N. J. Kong, J. J. Payne, J. Zhu and A. M. Johnson.
\emph{Saltation Matrices: The Essential Tool for Linearizing Hybrid
Dynamical Systems}. Proceedings of the IEEE 112(6), 585--608, 2024.
\href{https://arxiv.org/pdf/2306.06862}{Author manuscript, v3},
Section III-A/B; DOI 10.1109/JPROC.2024.3440211.
\bibitem{delle} M. L. Delle Monache, B. Piccoli and F. Rossi.
\emph{Traffic regulation via controlled speed limit}.
\href{https://arxiv.org/pdf/1603.04785}{arXiv:1603.04785}, 2016,
abstract and Section 3.3. Different control/objective; precedent only.
\end{thebibliography}
\end{document}
